\chapter{面向不确定输入的约束校准设计}
\section{总体方法}
原始观测先形成状态表征，再由校准模块生成监督对象。
\subsection{状态表征}
由于原始窗口仍含异步采样误差，后续校准不能直接使用原始观测。
为获得时间对齐能力，本节采用带掩码的状态编码器。
\begin{equation}
h_t = f(x_t, m_t)
\end{equation}
式中，$h_t$ 表示对齐后的状态，并作为校准模块的输入。
\subsection{监督校准}
由于状态表征尚未给出可信范围，校准器使用独立数据确定准入阈值。
该阈值把候选转换为带区间的监督对象，供下游预测器训练。
